모든 자료는 저장소의 \texttt{sphere\_packing\_research/} 아래에 있다(브랜치 \texttt{claude/laughing-allen-mobk2e}).
\begin{center}\small
\begin{longtable}{@{}p{5.6cm}p{10.4cm}@{}}
\toprule
경로 & 내용\\
\midrule
\endhead
\texttt{research\_log.md} & 회차별 연구 로그(부록 \ref{app:log} 와 동일)\\
\texttt{report/} & 이 보고서의 \LaTeX\ 원본과 \texttt{build.sh} (XeLaTeX + pandoc)\\
\texttt{rounds/R1\_lit\_*/notes.md} & 1회차 문헌 조사 노트 4종; \texttt{R1\_lit\_perfect/best\_known\_lattices.csv}\\
\texttt{rounds/R2\_klartag\_sim/} & Klartag 과정 시뮬레이터(\texttt{sim.py}), 결과표, 검증\\
\texttt{rounds/R2\_voronoi\_descent/} & Voronoi 형 탐색(\texttt{voronoi.py}), C 열거기 \texttt{fp\_enum.c}\\
\texttt{rounds/R3\_G\_klartag/}, \texttt{R3\_non\_free/}, \texttt{R3\_process/} & 이득 항등식, 텐서 안정성, no-go, Lemma N 의 증명 노트와 심사 보고\\
\texttt{rounds/R4\_perceptron\_*/} & 퍼셉트론 모형의 이론 노트와 LP 실험; \texttt{referee/} 에 $\chi^2$ 법칙의 증명\\
\texttt{rounds/R5\_leech13/} & Leech/$\Q(\sqrt{13})$ 의 정수 연산 계산(\texttt{stepA.py}--\texttt{stepF.py}), 심사자의 독립 재계산\\
\texttt{rounds/R5\_sl2p\_theory/} & 일반 $p$ 의 보조정리(A$'$, B, 직선 기준)와 심사\\
\texttt{rounds/R6\_quaternionic/} & \texttt{proof\_revised.md} (부록 \ref{app:proof}), 수치 검산 스크립트\\
\texttt{rounds/R7\_gluing/} & GGS 접착 데이터의 LP 전수 탐색(\texttt{glue\_search.py})과 심사\\
\texttt{rounds/R8\_equivariant/} & 등변 탐색(\texttt{eq.py}), 결과표, \texttt{best\_forms.json}, 검증\\
\texttt{rounds/R9\_trace\_lp\_lattices/} & 실제 격자의 대각합-LP(\texttt{tracelp.py}), 결과표, 검증\\
\texttt{rounds/R10\_weil\_lattices/} & $SL_2(p)$ Weil 격자: \texttt{orchestrator/} (\texttt{proto.py}, \texttt{run37.py}, \texttt{sub37.py}, \texttt{certify.py}, \texttt{roots37.py}), 빌더 \texttt{A/}, \texttt{B/}, 교차검증 \texttt{referee/}\\
\texttt{rounds/R11\_mordell\_weil/} & Mordell--Weil 재정식화 검증(T1), 반박 시도(T2)\\
\texttt{rounds/R12\_p61/} & $p=61$ 의 체계적 시험(\texttt{notes.md}, 정수 기저)과 독립 검증 \texttt{verify/}\\
\texttt{rounds/R10\_weil\_lattices/ orchestrator/scan.py} 등 & 13회차의 구조적 부분격자 스캔(\texttt{scan.py}, \texttt{rootsys.py}; 출력 \texttt{scan\_all\_output.txt}, 수정 후 재실행 \texttt{log\_scan3\_p*.txt})\\
\texttt{rounds/R13\_scan\_verify/} & 13회차 스캔의 독립 검증\\
\texttt{rounds/R14\_leech\_tensor/} & 비자유 작용의 틀 정리와 Leech 텐서족 정리(\texttt{proof.md}; 부록 \ref{app:proof14}), 문헌 점검(\texttt{novelty\_check.md}), 심사 보고 \texttt{referee\_A.md}, \texttt{referee\_B.md}, 수치 검산 스크립트\\
\texttt{rounds/R15\_rogers/} & $j$-틀 평균값 공식의 유도(\texttt{proof.md}; 부록 \ref{app:proof15})와 심사 보고 2건\\
\texttt{rounds/R16\_quotient\_curve/} & 몫곡선 사전(명제 \ref{prop:quotient})의 심사 보고\\
\texttt{rounds/R17\_gargava\_check/} & 인용한 Gargava 정리 16 의 $s\ge2$ 경우 증명 점검(\texttt{notes.md})과 심사 보고\\
\texttt{rounds/R18\_integration/} & 14·15회차 증명의 결합 점검과 정리 D·E 의 통합 진술(\texttt{report.md})\\
\texttt{rounds/R19\_classnumber/} & class number 가정을 제거하는 확장(\texttt{proof.md})과 심사 보고\\
\texttt{rounds/R20\_infinitely\_many\_stable/} & (기록 종료 후) ``안정한 $p$ 가 무한히 많은가'' 의 증명 시도: \texttt{proof.md} (부록 \ref{app:notes}), 정확한 ideal 격자 모형 \texttt{ideal\_model.py}, 전수 탐색 \texttt{scan\_exact.py} 와 로그, 다항식 모형 \texttt{mw\_poly\_model.py}, 버퍼를 키운 열거기 \texttt{fp\_enum\_big.c}; 저자의 독립 점검 \texttt{orchestrator\_check/}\\
\texttt{rounds/R21\_weaker\_conditions/} & (기록 종료 후) 가정의 약화 탐색: \texttt{notes.md} (부록 \ref{app:notes}), GAP 스캔 \texttt{scan.g}, \texttt{scan3.g}, \texttt{scan4.g} (마지막 것이 증명된 조건만 쓰는 최종판)와 출력 \texttt{scan\_out.txt}, \texttt{scan3\_lib.txt}, \texttt{scan4\_lib.txt}, $SL_2(q)$ 족 \texttt{sl2\_family.py}, 합성수 경우 \texttt{composite\_M.py}, \texttt{composite\_M\_cyclic.py}\\
\texttt{README.md} & 프로젝트 개요, 결과 요약, 재현 방법\\
\bottomrule
\end{longtable}
\end{center}
\paragraph{재현 예.} 구조적 부분격자 스캔: \texttt{cd rounds/R10\_weil\_lattices/orchestrator; python3 scan.py 37 61} ($p=37$ 은 수 초, $p=61$ 은 약 1분). $p=37$ 의 불안정성: \texttt{cd rounds/R10\_weil\_lattices/orchestrator; python3 run37.py 37 9,3; python3 sub37.py 37 3; python3 sub37.py 37 9 40; python3 certify.py 37 9 3 "[...]"} (마지막 인자는 \texttt{sub37.py} 가 출력하는 위반 벡터의 좌표). 필요한 것은 Python 3, numpy/scipy, gcc(\texttt{fp\_enum.c} 컴파일)뿐이다. 20회차의 정확한 모형: \texttt{cd rounds/R20\_infinitely\_many\_stable; python3 orchestrator\_check/chk.py 13 37 61} (랭크 24 격자의 근; 인자가 없으면 $p=13,37$; 수십 초, sympy 불필요). 21회차: \texttt{cd rounds/R21\_weaker\_conditions; python3 sl2\_family.py; python3 composite\_M.py} (sympy 필요), \texttt{gap -q -b scan4.g} (GAP 4.12 와 CTblLib 필요; 약 5분). \emph{주의}: \texttt{fp\_enum.c} (2·8회차, 10회차 빌더 A)는 좌표 버퍼가 64 개여서 랭크 64 이하에서만 써야 한다(\S\ref{sec:r20}); 그보다 큰 랭크에는 20회차 폴더의 \texttt{fp\_enum\_big.c} 또는 12회차 폴더의 \texttt{lat.c} 를 쓴다.
\paragraph{검증 수준에 대한 주의.} 이 보고서의 [증명됨] 은 ``하위 에이전트가 쓴 증명을 반박 담당 에이전트가 검토해 오류를 찾지 못했다'' 는 뜻이며 사람의 동료 심사를 거치지 않았다. 계산 결과 중 정수 연산으로 확정된 것([증명됨(계산기)])과 부동소수점 계산([수치관찰])을 구분해 표기했다. 인용 문헌 중 원문을 직접 확인하지 못한 것은 다음과 같다: Gross (1990), Nebe (1998), Dummigan, Tiep 의 $SL_2(p)$ 격자 논문들(호스트 접근 불가), Zassenhaus/Passman 의 Frobenius complement 구조정리와 Suzuki 의 분류(교과서적 사실로 기억에 의존해 인용), Weil (1946). 이들에 의존하는 진술에는 본문에서 ``인용'' 또는 ``원문 미열람'' 이라고 표시했다. 기록 종료 후의 20·21회차(\S\ref{sec:r20}, \S\ref{sec:r21})는 증명에 대한 정식 AI 심사를 받지 않았다(서술을 원자료와 대조하는 비평만 한 차례 거쳤다). arXiv 논문(Klartag, Rezvyakova, AGZ, GGS, Gargava, Kim 등)은 본문 또는 초록을 직접 열람했으나 일부는 요약 도구를 거쳐 읽었다.
